\paragraph{Chebotarev 商熵、可见角色能量与平方根统计障碍}
在 primitive 同余层的严格 Fourier 反演、subgroup--Möbius 完全反演、商映射下的精确能量损失、非阿贝尔通道能量账本、最小偏差承载商以及主导通道见证集合都已建立之后，还可把群商粗化、Plancherel 通道分裂、次主奇点与 RH/GRH 型统计后果进一步压缩为下列四条结构后果。它们分别把商映射的 \(L^2\) 收缩提升为相对熵链式律，把可见偏差写成可见/隐藏角色能量的精确正交分解，把主导扭曲通道的次主解析障碍严格识别为表示维数加权极点，并把平方根谱隙条件转译为 TV/\(\chi^2\)/KL 的统一指数律及其见证性反证。

\begin{theorem}[群商下 Chebotarev 分布相对均匀基线的 KL 链式律]\label{thm:xi-time-part64ac-abelian-quotient-kl-chain}
沿用推论 \ref{cor:kernel-chebotarev-L2-monotone} 的记号。设
\[
\pi:G_2\twoheadrightarrow G_1
\]
为有限阿贝尔群满同态，并记 \(K:=\ker\pi\)。对固定 \(u\in(0,1]\) 与 \(n\ge 1\)，定义两层 primitive Chebotarev 概率分布
\[
\mu_{n,u}^{(i)}(c_i):=\frac{B_{n,c_i}^{(i)}(u)}{B_n(u)},
\qquad
B_n(u):=B_n^{(1)}(\mathbf{1};u)=B_n^{(2)}(\mathbf{1};u),
\qquad
i\in\{1,2\},
\]
并记 \(u_{G_i}\) 为 \(G_i\) 上的均匀分布。对每个满足 \(\mu_{n,u}^{(1)}(c_1)>0\) 的 \(c_1\in G_1\)，定义纤维条件分布
\[
\mu_{n,u}^{(c_1),\mathrm{fib}}(c_2)
:=
\frac{\mu_{n,u}^{(2)}(c_2)}{\mu_{n,u}^{(1)}(c_1)},
\qquad
c_2\in \pi^{-1}(c_1),
\]
并记 \(u_{\pi^{-1}(c_1)}\) 为该纤维上的均匀分布。则有严格链式恒等式
\[
\boxed{\;
D_{\mathrm{KL}}\!\bigl(\mu_{n,u}^{(2)}\Vert u_{G_2}\bigr)
=
D_{\mathrm{KL}}\!\bigl(\mu_{n,u}^{(1)}\Vert u_{G_1}\bigr)
+
\sum_{c_1\in G_1}\mu_{n,u}^{(1)}(c_1)\,
D_{\mathrm{KL}}\!\bigl(\mu_{n,u}^{(c_1),\mathrm{fib}}\Vert u_{\pi^{-1}(c_1)}\bigr)\; }.
\]
其中对 \(\mu_{n,u}^{(1)}(c_1)=0\) 的纤维约定对应项为 \(0\)。特别地，商映射所丢失的全部相对熵恰等于纤维条件非均匀性的平均 KL 成本。
\leanverified{paper\_xi\_time\_part64ac\_abelian\_quotient\_kl\_chain}
\end{theorem}
\begin{proof}
由于 \(\mu_{n,u}^{(1)}=\pi_\ast\mu_{n,u}^{(2)}\) 且 \(\pi\) 为群满同态，每条纤维大小都等于 \(|K|\)，故
\[
u_{G_2}(c_2)
=
\frac{1}{|G_2|}
=
\frac{1}{|G_1|\,|K|}
=
u_{G_1}(\pi(c_2))\,u_{\pi^{-1}(\pi(c_2))}(c_2).
\]
于是对任意满足 \(\mu_{n,u}^{(2)}(c_2)>0\) 的 \(c_2\in G_2\)，有
\[
\log\frac{\mu_{n,u}^{(2)}(c_2)}{u_{G_2}(c_2)}
=
\log\frac{\mu_{n,u}^{(1)}(\pi(c_2))}{u_{G_1}(\pi(c_2))}
+
\log\frac{\mu_{n,u}^{(\pi(c_2)),\mathrm{fib}}(c_2)}{u_{\pi^{-1}(\pi(c_2))}(c_2)}.
\]
对 \(\mu_{n,u}^{(2)}\) 取期望并按纤维分组，即得所示分解；这也是引理 \ref{lem:pom-fiberwise-kl} 在等势群商情形下的直接专门化。证毕。
\end{proof}

\begin{theorem}[可见/隐藏角色能量的精确正交分解]\label{thm:xi-time-part64ac-visible-hidden-character-energy-splitting}
沿用定理 \ref{thm:xi-time-part64ac-abelian-quotient-kl-chain} 的设定，并记
\[
r_n^{(i)}(c_i;u):=|G_i|\,\mu_{n,u}^{(i)}(c_i)
=
\frac{|G_i|\,B_{n,c_i}^{(i)}(u)}{B_n(u)},
\qquad i\in\{1,2\}.
\]
把 \(\widehat{G_1}\) 通过 \(\pi^\ast\) 识别为
\[
K^\perp:=\{\chi\in \widehat{G_2}:\ \chi|_K\equiv 1\}\subseteq \widehat{G_2},
\]
并记细层角色幅度
\[
B_n^{(2)}(\chi;u):=\sum_{c_2\in G_2}\chi(c_2)\,B_{n,c_2}^{(2)}(u).
\]
则有两组精确分解：
\[
\boxed{\;
\bigl\|r_n^{(2)}(\cdot;u)-1\bigr\|_{L^2(G_2)}^2
=
\bigl\|r_n^{(1)}(\cdot;u)-1\bigr\|_{L^2(G_1)}^2
+
\frac{1}{|G_2|}
\sum_{c_1\in G_1}\sum_{c_2\in\pi^{-1}(c_1)}
\bigl|r_n^{(2)}(c_2;u)-r_n^{(1)}(c_1;u)\bigr|^2\; }.
\]
\[
\boxed{\;
\bigl\|r_n^{(1)}(\cdot;u)-1\bigr\|_{L^2(G_1)}^2
=
\frac{1}{B_n(u)^2}
\sum_{\chi\in K^\perp\setminus\{\mathbf{1}\}}
\bigl|B_n^{(2)}(\chi;u)\bigr|^2\; }.
\]
\[
\boxed{\;
\bigl\|r_n^{(2)}(\cdot;u)-1\bigr\|_{L^2(G_2)}^2
-
\bigl\|r_n^{(1)}(\cdot;u)-1\bigr\|_{L^2(G_1)}^2
=
\frac{1}{B_n(u)^2}
\sum_{\chi\in \widehat{G_2}\setminus K^\perp}
\bigl|B_n^{(2)}(\chi;u)\bigr|^2\; }.
\]
因而，群商仅保留在 \(K\) 上平凡的可见角色通道，并精确删除全部隐藏通道能量。
\end{theorem}
\begin{proof}
由推论 \ref{cor:kernel-chebotarev-L2-monotone} 的证明，
\[
r_n^{(1)}(\cdot;u)\circ\pi
=
\mathbb{E}_\pi\!\bigl[r_n^{(2)}(\cdot;u)\bigr]
\]
正是 \(r_n^{(2)}(\cdot;u)\) 沿纤维常值子空间的条件期望。故在 \(L^2(G_2)\) 中可应用正交投影的勾股恒等式：
\[
\bigl\|r_n^{(2)}-1\bigr\|_2^2
=
\bigl\|r_n^{(1)}\circ\pi-1\bigr\|_2^2
+
\bigl\|r_n^{(2)}-r_n^{(1)}\circ\pi\bigr\|_2^2.
\]
按纤维展开第二项，即得第一式。

另一方面，推论 \ref{cor:kernel-chebotarev-plancherel-energy} 应用于细层给出
\[
\bigl\|r_n^{(2)}(\cdot;u)-1\bigr\|_{L^2(G_2)}^2
=
\frac{1}{B_n(u)^2}
\sum_{\chi\in \widehat{G_2}\setminus\{\mathbf{1}\}}
\bigl|B_n^{(2)}(\chi;u)\bigr|^2.
\]
而粗层角色正好对应于 \(K^\perp=\pi^\ast(\widehat{G_1})\)，故同理
\[
\bigl\|r_n^{(1)}(\cdot;u)-1\bigr\|_{L^2(G_1)}^2
=
\frac{1}{B_n(u)^2}
\sum_{\chi\in K^\perp\setminus\{\mathbf{1}\}}
\bigl|B_n^{(2)}(\chi;u)\bigr|^2,
\]
得到第二式。将两式相减即得第三式；这也是定理
\ref{thm:xi-time-part63b-primitive-quotient-energy-loss-exact}
在相对 Haar 密度口径下的直接重写。证毕。
\end{proof}

\begin{theorem}[主导非平凡通道的极点阶与见证振荡]\label{thm:xi-time-part64ac-dominant-channel-pole-order-witness-oscillation}
沿用定理 \ref{thm:xi-time-part64-cover-zeta-pole-order-weighted-channel-multiplicity} 与定理 \ref{thm:kernel-chebotarev-second-main-term-witness} 的记号。固定 \(u\in(0,1]\)，并假设存在唯一非平凡通道 \(\rho_\star\neq \mathbf{1}\) 与其简单本征值 \(\mu_\star\in\Spec(M_{K,\rho_\star}(u))\) 满足
\[
|\mu_\star|
=
\max_{\rho\neq \mathbf{1}}\rho\!\bigl(M_{K,\rho}(u)\bigr)
=:\eta(u),
\]
且该通道在 \(\mu_\star\) 处具有谱隙。则有：
\begin{enumerate}
  \item 去主因子覆盖 \(\zeta\) 在 \(z=\mu_\star^{-1}\) 处具有精确极点阶
  \[
  \boxed{\;
  \operatorname{ord}_{z=\mu_\star^{-1}}
  \frac{\zeta_{K,G}(z,u)}{L_{K,\mathbf{1}}(z,u)}
  =
  \dim\rho_\star\; }.
  \]
  \item 存在某个常数 \(\eta_-(u)<\eta(u)\)，使得对任意共轭不变集合 \(A\subseteq G\)，若其指示函数 \(1_A\) 在 \(\rho_\star\)-通道上的 Fourier 分量非零，则记
  \[
  B_{n,A}(u):=\sum_{C\subseteq A}B_{n,C}(u),
  \qquad
  \mu_{\mathrm{Haar}}(A):=\frac{|A|}{|G|},
  \]
  并存在 \(c_A\neq 0\) 使
  \[
  \boxed{\;
  B_{n,A}(u)
  =
  \mu_{\mathrm{Haar}}(A)\frac{\lambda_1(u)^n}{n}
  +
  \Re\!\bigl(c_A\mu_\star^n\bigr)\frac1n
  +
  O\!\left(\frac{\eta_-(u)^n}{n}\right)\; }.
  \]
  特别地，存在常数 \(c(A)>0\) 使
  \[
  \boxed{\;
  \left|B_{n,A}(u)-\mu_{\mathrm{Haar}}(A)\frac{\lambda_1(u)^n}{n}\right|
  \ge
  c(A)\frac{\eta(u)^n}{n}\; }
  \]
  在无穷多个 \(n\) 上成立。
\end{enumerate}
\end{theorem}
\begin{proof}
由 Artin 因子化
\[
\frac{\zeta_{K,G}(z,u)}{L_{K,\mathbf{1}}(z,u)}
=
\prod_{\rho\neq \mathbf{1}}L_{K,\rho}(z,u)^{\dim\rho}
\]
可知，该商在 \(z=\mu_\star^{-1}\) 处的极点阶只由非平凡通道决定。再由定理 \ref{thm:xi-time-part64-cover-zeta-pole-order-weighted-channel-multiplicity}，
\[
\operatorname{ord}_{z=\mu_\star^{-1}}
\frac{\zeta_{K,G}(z,u)}{L_{K,\mathbf{1}}(z,u)}
=
\sum_{\rho\in\widehat G\setminus\{\mathbf{1}\}}(\dim\rho)\,m_\rho(\mu_\star),
\]
其中 \(m_\rho(\mu_\star)\) 为 \(\mu_\star\) 在 \(M_{K,\rho}(u)\) 中的代数重数。由唯一主导通道假设与 \(\mu_\star\) 的简单性可知，\(\rho_\star\) 贡献 \(m_{\rho_\star}(\mu_\star)=1\)，而其余非平凡通道都不含该特征值，故极点阶等于 \(\dim\rho_\star\)。

第二部分正是定理 \ref{thm:kernel-chebotarev-second-main-term-witness} 的直接应用：唯一主导通道给出单一的次主振荡模态，谱隙使所有其余通道并入某个 \(\eta_-(u)<\eta(u)\) 的统一误差项；若 \(1_A\) 在 \(\rho_\star\)-通道上的分量非零，则对应系数 \(c_A\neq 0\)，从而沿无穷多个 \(n\) 得到 \(\Omega(\eta(u)^n/n)\) 的见证振荡下界。证毕。
\end{proof}

\begin{theorem}[GRH 型平方根谱隙的 TV/\texorpdfstring{$\chi^2$}{chi2}/KL 统计后果]\label{thm:xi-time-part64ac-grh-tv-chi2-kl-barrier}
沿用定理 \ref{thm:xi-time-part64ab-primitive-chebotarev-haar-tv-exponent} 的设定，并进一步假设 \(G\) 为有限阿贝尔群。定义
\[
\mu_{n,u}(c):=\frac{B_{n,c}(u)}{B_n(u)},
\qquad
u_G(c):=\frac{1}{|G|},
\qquad
r_n(c;u):=|G|\,\mu_{n,u}(c),
\]
以及
\[
\eta(u):=\max_{\chi\neq \mathbf{1}}\rho\!\bigl(M_{K,\chi}(u)\bigr).
\]
则有统一指数界
\[
\boxed{\;
\bigl\|\mu_{n,u}-u_G\bigr\|_{\mathrm{TV}}
=
O\!\left(\Bigl(\frac{\eta(u)}{\lambda_1(u)}\Bigr)^n\right),\;}
\]
\[
\boxed{\;
\chi^2\!\bigl(\mu_{n,u}\Vert u_G\bigr)
=
O\!\left(\Bigl(\frac{\eta(u)}{\lambda_1(u)}\Bigr)^{2n}\right),\qquad
D_{\mathrm{KL}}\!\bigl(\mu_{n,u}\Vert u_G\bigr)
=
O\!\left(\Bigl(\frac{\eta(u)}{\lambda_1(u)}\Bigr)^{2n}\right).\;}
\]
特别地，若 \(\textsf{GRH}_{K,G}(u)\) 成立，即
\[
\eta(u)\le \lambda_1(u)^{1/2},
\]
则
\[
\boxed{\;
\bigl\|\mu_{n,u}-u_G\bigr\|_{\mathrm{TV}}=O\!\bigl(\lambda_1(u)^{-n/2}\bigr),\qquad
\chi^2\!\bigl(\mu_{n,u}\Vert u_G\bigr),\
D_{\mathrm{KL}}\!\bigl(\mu_{n,u}\Vert u_G\bigr)
=
O\!\bigl(\lambda_1(u)^{-n}\bigr).\;}
\]
反之，若定理 \ref{thm:xi-time-part64ac-dominant-channel-pole-order-witness-oscillation} 的唯一主导通道假设成立，且
\[
\eta(u)>\lambda_1(u)^{1/2},
\]
则存在某个集合 \(A\subseteq G\) 与常数 \(c(A)>0\)，使
\[
\boxed{\;
\bigl|\mu_{n,u}(A)-u_G(A)\bigr|
\ge
c(A)\Bigl(\frac{\eta(u)}{\lambda_1(u)}\Bigr)^n\; }
\]
在无穷多个 \(n\) 上成立。因而此时不可能存在统一的平方根级 TV 收敛律。
\end{theorem}
\begin{proof}
阿贝尔情形下，共轭类分布即元素分布，故定理 \ref{thm:xi-time-part64ab-primitive-chebotarev-haar-tv-exponent} 直接给出 TV 上界。另一方面，
\[
\chi^2\!\bigl(\mu_{n,u}\Vert u_G\bigr)
=
\frac{1}{|G|}\sum_{c\in G}\bigl(r_n(c;u)-1\bigr)^2.
\]
由推论 \ref{cor:kernel-chebotarev-plancherel-energy} 的相对 Haar 密度版本，
\[
\chi^2\!\bigl(\mu_{n,u}\Vert u_G\bigr)
=
\frac{1}{B_n(u)^2}
\sum_{\chi\in\widehat G\setminus\{\mathbf{1}\}}
\bigl|B_n(\chi;u)\bigr|^2.
\]
对每个非平凡角色 \(\chi\neq \mathbf{1}\)，由命题 \ref{prop:kernel-chebotarev-mobius-adams} 与谱半径估计有
\[
B_n(\chi;u)=O\!\left(\frac{\eta(u)^n}{n}\right),
\]
而 \(B_n(u)\sim \lambda_1(u)^n/n\) 由定理 \ref{thm:kernel-weighted-prime-orbit} 给出。由于非平凡角色数有限，遂得
\[
\chi^2\!\bigl(\mu_{n,u}\Vert u_G\bigr)
=
O\!\left(\Bigl(\frac{\eta(u)}{\lambda_1(u)}\Bigr)^{2n}\right).
\]
再由标准不等式
\[
D_{\mathrm{KL}}(P\Vert Q)\le \chi^2(P\Vert Q),
\]
立即得到 KL 上界。若 \(\eta(u)\le \lambda_1(u)^{1/2}\)，则把该平方根界代入上述三式即得所示结论。

反向方向中，若唯一主导通道存在且 \(\eta(u)>\lambda_1(u)^{1/2}\)，则由定理
\ref{thm:xi-time-part64ac-dominant-channel-pole-order-witness-oscillation}
可取某个 \(A\subseteq G\) 使
\[
\left|B_{n,A}(u)-u_G(A)\frac{\lambda_1(u)^n}{n}\right|
\ge
c(A)\frac{\eta(u)^n}{n}
\]
在无穷多个 \(n\) 上成立。再除以
\[
B_n(u)\sim \frac{\lambda_1(u)^n}{n}
\]
即得
\[
\bigl|\mu_{n,u}(A)-u_G(A)\bigr|
\ge
c'(A)\Bigl(\frac{\eta(u)}{\lambda_1(u)}\Bigr)^n
\]
沿无穷多个 \(n\) 成立。由于总变差是事件偏差的上确界，若再有统一的 \(O(\lambda_1(u)^{-n/2})\) 上界，则会与 \(\eta(u)>\lambda_1(u)^{1/2}\) 矛盾。证毕。
\end{proof}

\noindent
因此，primitive Chebotarev 的 quotient pushforward 不仅满足 \(L^2\) 收缩，而且在相对熵层面具有严格链式记账；相应的可见偏差可精确拆分为可见/隐藏角色通道能量；一旦某个非平凡通道主导全部次主谱，它便同时决定去主因子覆盖 \(\zeta\) 的最近奇点阶数与 Chebotarev 见证集合的振荡量级；而 RH/GRH 型平方根谱隙条件则恰把这些通道参数进一步压缩为 TV/\(\chi^2\)/KL 的统一统计指数与尖锐障碍。于是，商粗化、角色 Fourier、Artin 因子化与平方根误差预算便在同一条有限维通道链上闭合。

\endinput
